\section{从预测到净收益的统计套利研究}
\label{ch:lower-stat-arb-research}

\begin{itemize}
  % Generated from curriculum/manifest.json; do not edit by hand.
\item 直接先修：下册《估计与状态推断：滤波、切换与变点》。

\end{itemize}

均值回归预测要先变成两腿订单，才能计算损益。以下区分预测变化、目标仓位和实际成交，并说明对冲比率更新为什么不等于持仓盈利。


\MFQTransition{先画时间轴，再拟合模型}

标签跨越 $h$ 期时，应删除结果尚未在拟合截止前完成、或与验证标签共享信息的训练样本。额外缓冲由标签完成、发布延迟与依赖结构决定，不是任何验证后都必须固定空置若干期。事先确定示例为训练索引 $0{:}5$、验证 $8{:}9$、交易 $12{:}13$，$h=2$、embargo=2。若训练扩到 6，其标签触及验证起点，验证器立即拒绝。随机 K 折、全样本 scaler、用交易期重新选阈值也由同一信息时间轴判断，而不是靠函数名称判断。

% Generated by tools/build_figures.py; edit figures/figure-specs.json instead.
\MFQFigure{tex/figures/generated/lower-sa-timeline.pdf}{滚动研究先用已有数据估计关系及选择阈值，再观察后续交易区间；重叠标签需要单独处理。}{lower-sa-timeline}


\MFQTransition{预测必须真实决定仓位}

本节先把预测明确定义为下一期价差变化 $\widehat{\Delta z}_{t+1|t}$，而非下一期价差水平。预测变化为正时做多价差，决策映射为
\[
s_t=\begin{cases}
+1,&\widehat{\Delta z}_{t+1|t}\ge c,\\
-1,&\widehat{\Delta z}_{t+1|t}\le-c,\\
0,&|\widehat{\Delta z}_{t+1|t}|<c,
\end{cases}
\qquad w_t=\min(L,|s_t|L)\operatorname{sign}(s_t).
\]
同时声明阈值 $c=0.5$、上限 $L=1$、持有期 1、每期调仓。预测 $[0.8,-0.7,0.1,0.9]$ 因而产生目标仓位 $[1,-1,0,1]$，不是从 解析参照 直接读取交易收益。

% Generated by tools/build_figures.py; edit figures/figure-specs.json instead.
\MFQFigure{tex/figures/generated/lower-sa-position-pnl.pdf}{预测先映射为有界仓位，仓位再作用于下一期收益；换手成本在仓位变化时扣除。}{lower-sa-position-pnl}


成交比例作用于订单增量，而不是目标仓位本身。令 $p_{t-1}$ 为旧仓位、$f_t$ 为成交比例，则
\[
q_t=w_t-p_{t-1},\qquad p_t=p_{t-1}+f_tq_t.
\]
成交比例 $[1,0.5,1,0]$ 因而产生实际仓位 $[1,0,0,0]$：第二期从 $+1$ 翻到 $-1$ 的订单量是 $-2$，成交一半后恰好平仓。最后一笔未成交时，仓位保持上一期的 0；更一般地，零成交必须保持旧仓位，而非强制归零。毛收益逐期为 $p_t r_{t+1}$，换手从空仓开始：
\[
\text{turnover}=|1-0|+|0-1|+|0-0|+|0-0|=2.
\]
实现收益 $[0.02,-0.01,0.03,0.04]$ 给出毛收益 $0.02$；单位换手成本 $0.002$ 给出成本 $0.004$，故 $R^{\mathrm{net}}=0.016$。把成本率加倍后，净收益应再下降 $0.004$，这是计算记录的敏感性检查。

透明实现逐期暴露订单、成交与旧仓位的语义；成熟库对照把同一递推写成有限维下三角线性系统并交给 SciPy 求解。后者适合独立核对小型事先确定计算记录，但稠密矩阵会随期数平方增长，也没有队列位置、撮合优先级、冲击或停牌状态，因此不能冒充真实执行引擎。长时间轴应利用双对角稀疏结构或直接递推，市场微观结构则需要另建状态模型。

\MFQTransition{失效检测与研究材料}

研究材料必须逐窗保存：估计截止、模型参数、残差诊断、状态概率、预测、目标仓位、成交比例、实际仓位、换手、成本与净收益。失效规则可包括协整残差单位根恶化、半衰期超出持有期、状态概率熵升高、变点警报和成交率下降。规则必须在交易期前事先确定；用全样本挑选“最早正确警报”没有意义。

错误研究材料审计要追问：预测是否真实进入仓位？收益日期是否晚于决策？未成交是否仍计收益？成本是否按实际仓位变化计算？scaler、协整向量与阈值是否只在训练窗口拟合？这些问题比多加一个模型名称更能决定研究可信度。

若数据不支持交易时点，项目应停在方法演示，不得包装成回测。

\MFQTransition{从价差 z-score 到目标仓位}

训练窗口估计价差均值和波动后，$z_t=(s_t-\hat\mu_t)/\hat\sigma_t$。最简单规则在 $z_t>z_{\rm entry}$ 时做空价差、$z_t<-z_{\rm entry}$ 时做多，在 $|z_t|<z_{\rm exit}$ 时平仓。入口和出口不同形成 hysteresis，减少阈值附近反复交易。

仓位可固定，也可按预期回归幅度、风险和半衰期缩放。无论哪种，必须设最大持有期、止损、关系失效和不可交易处理。参数只在内层验证选择；用全样本最优 z-score 是数据挖掘。

\MFQTransition{配对交易的自融资计算记录}

若价差为 $s_t=y_t-\beta_tx_t$，做空一单位价差意味着卖出 $y$、买入 $\beta_t$ 单位 $x$。beta 变化会触发再平衡成本。收益应从两腿实际持仓与下一期价格变化计算：
\[
\Pi_{t+1}-\Pi_t=w_{y,t}\Delta y_{t+1}
+w_{x,t}\Delta x_{t+1}-C_{t+1}.
\]
直接用 $-\Delta s$ 当收益只在 beta 固定、单位一致且忽略融资/成本时成立。

\MFQTransition{候选选择与多重检验}

从 1000 只股票枚举约 50 万对，再挑最平稳价差，会产生巨大选择偏差。先用行业、供应链或共享经济暴露缩小候选，再在训练期做协整检验；完整记录尝试数，并在独立未来期评价整个筛选程序。

篮子套利用 Johansen 或稀疏方法寻找多资产组合，实施复杂度随腿数上升：更多点差、借券和异步成交。统计上更平稳的篮子不一定更可交易。

\MFQTransition{收益分解与压力测试}

将净收益拆成均值回归收益、beta 更新、融资、点差、冲击、借券和未成交。压力测试包括相关关系突然断裂、波动翻倍、半衰期延长、单腿停牌和借券召回。失效规则应在压力前定义，而非根据损失路径挑一个恰好提前退出的阈值。

\begin{diagnostic}
如果预测序列与策略持仓没有确定映射，或收益来自 合成样本 中预先给定的 \texttt{trade\_return}，那只是并排展示模型和交易数字，不是模型驱动策略。
\end{diagnostic}

\section{本路线的综合项目}

选定一条协整或状态切换关系，沿时间轴完成候选选择、滚动估计、仓位映射和成本后评价。综合报告应同时展示关系失效、参数漂移、未成交和借券约束的一个反例，并把结论限定在声明的资产池与研究期内。

\subsection{变动 beta 为什么产生虚假的价差收益}

设 $x_t=50,y_t=101,\beta_t=2$，下一期 $x_{t+1}=51,y_{t+1}=102$。做空一单位价差的持仓为 $w_y=-1,w_x=2$，不计成本的损益为 $-1(102-101)+2(51-50)=1$。它等于固定 beta 下的 $-\Delta(y-2x)$。

若下一期重新估计得到 $\beta_{t+1}=2.1$，新定义价差为 $102-2.1(51)=-5.1$，相较旧价差 1 下降 6.1，却不是本期赚了 6.1。恒等式为
\[
 \Delta(y-\beta x)=\Delta y-\beta_t\Delta x-x_{t+1}\Delta\beta.
\]
最后一项 $-51(0.1)=-5.1$ 是定义更新，不是旧持仓损益。要把持仓调到新的 beta，还需买入 $0.1$ 单位 $x$，并在现金账户记录支出及成本。notebook 将直接显示两腿股数、价格变化与现金调整，避免用重新拟合的历史价差替代真实交易。

\section{综合练习}
\begin{enumerate}[leftmargin=*]
  \item \textbf{口述概念：}为什么 ADF 的 t 统计量不能查普通 t 分布？
  \item \textbf{概念：}高相关与协整分别约束什么？
  \item \textbf{笔试推导：}由离散 AR 系数 $\phi$ 推导 OU 的 $\kappa$ 与半衰期。
  \item \textbf{推导：}从 VAR 写出 VECM，并解释 $\Pi=\alpha\beta^{\mathsf T}$。
  \item \textbf{数值编程：}改变残差 $\phi$，画半衰期敏感性并标记无效区间。
  \item \textbf{编程：}比较透明 DF、\texttt{adfuller} 与 Engle--Granger 专用临界值。
  \item \textbf{研究判断：}一份报告先从 500 对资产中选最平稳价差，再对所选对报告 p 值。指出选择偏差并设计修复。
\end{enumerate}

\subsection*{基础衔接：独立计算}
$\phi=0.5$，采样间隔两天，求 OU 均值偏离的半衰期。

\subsection*{模拟频率与尺度}
\begin{enumerate}[leftmargin=*]
\item 共同趋势加平稳价差与独立随机游走，两组拒绝频率各估计什么？拒绝率 0.5 时，60 次和 300 次重复的模拟标准误各为多少？
\item 固定 $D,z$，将 $y=D\beta+cz$ 的 $c$ 从 0.5 改为 2，推导系数误差、残差与固定规格 DF 统计量的变化。
\end{enumerate}

\section{分层练习}
\begin{enumerate}[leftmargin=*]
  \item \textbf{口述概念：}写出预测、滤波、平滑各自条件的信息集。
  \item \textbf{概念：}为什么 regime 标签会交换？
  \item \textbf{笔试推导：}推导标量 Kalman 增益与后验方差。
  \item \textbf{推导：}写出两状态过滤的一步预测与 Bayes 更新。
  \item \textbf{数值编程：}改变 $Q/R$，记录滤波路径的响应速度。
  \item \textbf{编程：}用多组初值拟合 Markov switching，比较似然与状态解释。
  \item \textbf{研究判断：}报告用全样本平滑状态回填历史交易。构造能执行的负例。
\end{enumerate}

\MFQTransition{仓位与收益计算与应用分析题组}
\begin{enumerate}[leftmargin=*]
  \item \textbf{口述概念：}purge 与 embargo 分别阻断什么信息路径？
  \item \textbf{概念：}为什么目标仓位不能直接用于收益计算？
  \item \textbf{笔试推导：}逐期重建本章 $0.02$ 毛收益、2 单位换手和 $0.016$ 净收益。
  \item \textbf{数值编程：}把持有期改为 2，解释仓位和换手变化。
  \item \textbf{编程：}加入借券失败与停牌成交比例，扩展失败状态而非删除样本。
  \item \textbf{研究判断：}一项研究用完整样本拟合标准化参数，并按未成交的目标仓位计算收益。分别说明两种做法改变了什么信息或持仓假设，并重新写出正确的计算顺序。
  \item \textbf{综合项目：}选择一组可用数据，估计价差并制定仓位规则。比较简单基线，计算部分成交和成本后的结果，再分析关系变化时会发生什么。
\end{enumerate}

\subsection*{基础衔接：独立计算}
若上述交易的两腿总成本为 0.2，净损益是多少？下一期 beta 改变是否改变本期已实现损益？

